\subsection{Dini moduli and limits of distances}
\label{sec:dini-metric-distance}

The annular argument of \cref{thm:holder-current-distance} also applies to a continuous potential whose modulus of continuity is not a power; what matters is a Dini condition on the square root of the modulus. This gives uniform distance estimates for families whose H\"older constants are not controlled.

Keep the coordinate domain $U\subset\C^n$ with its Euclidean form $\omega_{\mathrm E}$ and volume $\mathrm{d}\lambda$, but now let $v\in C^0(U)$ and let $\theta$ be a smooth real $(1,1)$-form with $T=\theta+\ddc v\geq0$. On the relatively compact charts under consideration, choose a nondecreasing function $m\colon(0,r_0)\to[0,+\infty)$ with $m(r)\to0$ as $r\to0$ such that
\[
 |v(x)-v(y)|\leq m(|x-y|),\qquad |x-y|<r_0.
\]
Put
\begin{equation}\label{eq:dini-distance-modulus}
 D(r)=r+\int_0^r\frac{\sqrt{m(t)}}{t}\,\mathrm{d}t.
\end{equation}

\Needspace{14\baselineskip}
\begin{theorem}\label{thm:dini-current-distance}\index{Dini modulus!metric distance}\index{Metric completion!Dini potential}
Assume that $D(r_0/2)<+\infty$ and that $T$ is a smooth Hermitian form $\eta$ on $U\setminus A$ for a proper analytic subset $A\subset U$. Then every point of $U$ has a neighborhood $V\Subset U$ with constants $C,r_1>0$ such that, for $x,y\in V$,
\[
 d_\eta^U(x,y)\leq C D(8|x-y|),
 \qquad x,y\notin A,\quad |x-y|<r_1.
\]
The constants depend only on the charts and an upper bound for $|\theta|$; the dependence on $v$ enters only through $D$. All sequences in $V\setminus A$ converging to a given point of $A\cap V$ have a common limit in the metric completion. If these hypotheses hold on a finite cover of a compact connected manifold by charts, the metric completion of the regular locus is compact, and the identity of the regular locus extends to a continuous surjection of the manifold onto it. If $A$ is nonempty, the metric $\eta$ on the regular locus is incomplete.
\end{theorem}
\begin{proof}
Write $d=2n$ and $\mu=T\wedge\omega_{\mathrm E}^{n-1}$. The same cutoff calculation as in \cref{thm:holder-current-distance}, with a constant subtracted from $v$, gives
\[
 \mu(B(x,s))\leq C\bigl(s^d+s^{d-2}m(2s)\bigr).
\]
This uses only the bound for $\theta$; closedness of $\theta$ is not needed. If $\tau$ is the Euclidean trace of $\eta$, then $\tau\,\mathrm{d}\lambda\leq C\mu$ on the regular locus. Polar integration bounds the average length of the segments from a regular point $x$ to the points of $B(x,R)$ by
\[
 C\int_{B(x,R)}\sqrt{\tau(\zeta)}
                      |\zeta-x|^{1-d}\,\mathrm{d}\lambda(\zeta).
\]
On an annulus of outer radius $s$, the Cauchy--Schwarz inequality and the trace bound make the contribution to this integral at most
\[
 Cs^{1-d}s^{d/2}
       \bigl(s^d+s^{d-2}m(2s)\bigr)^{1/2}
 \leq C\bigl(s+\sqrt{m(2s)}\bigr).
\]
For $s_j=2^{-j}R$, monotonicity gives
\[
 \sum_{j\geq0}\sqrt{m(2s_j)}
 \leq\frac1{\log2}\int_0^{4R}\frac{\sqrt{m(t)}}t\,\mathrm{d}t.
\]
Indeed, $\sqrt{m(2s_j)}\log2$ is at most the integral of $\sqrt{m(t)}/t$ over $[2s_j,4s_j]$, and these intervals have disjoint interiors; no doubling condition on $m$ is needed. The average length is therefore at most $CD(4R)$.

For $x,y\in V\setminus A$, average the lengths of the broken segments $[x,z]\cup[z,y]$ over $z$ in the ball of radius $|x-y|/2$ about the midpoint of $x$ and $y$, applying the preceding bound with $R=2|x-y|$ at each endpoint. Since the volume of this ball is comparable to $R^d$, some broken path has length at most $CD(8|x-y|)$. As in the proof of \cref{thm:holder-current-distance}, almost every such path avoids $A$, and for small $V$ it stays inside $U$.

Since $D(r_0/2)<+\infty$, we have $D(r)\to0$ as $r\to0$. Hence any two regular sequences converging to the same point are Cauchy and have the same limit in the completion. The argument for the completion in the proof of \cref{thm:holder-current-distance} applies with this modulus. A positive lower bound for $\eta$ near a regular point prevents a sequence converging to a point of $A$ from converging in the metric to the regular point. This also proves incompleteness when $A$ is nonempty.
\end{proof}

\Needspace{10\baselineskip}
\begin{corollary}\label{cor:dini-distance-rates}\index{Dini modulus!logarithmic distance rate}
In the setting of \cref{thm:dini-current-distance}, let $V$ be as in the theorem. For $0<\beta\leq1$, a modulus $m(r)\leq Cr^\beta$ gives $d_\eta^U(x,y)\leq C'|x-y|^{\beta/2}$ for sufficiently close $x,y\in V\setminus A$. If instead
\[
 m(r)\leq C(-\log r)^{-a},\qquad a>2,
\]
then for sufficiently close $x,y\in V\setminus A$,
\[
 d_\eta^U(x,y)\leq C'(-\log|x-y|)^{1-a/2}.
\]
These are only sufficient conditions; we do not claim that the Dini condition is necessary.
\end{corollary}
\begin{proof}
Integrating the square root of the power modulus gives a multiple of $r^{\beta/2}$. For the logarithmic modulus, the substitution $s=-\log t$ gives
\[
 \int_0^r\frac{(-\log t)^{-a/2}}t\,\mathrm{d}t
 =\frac{(-\log r)^{1-a/2}}{a/2-1}.
\]
In either case the term $r$ in $D(r)$ is bounded by the asserted rate for small $r$. The theorem now gives both estimates; the fixed dilation of the argument of $D$ changes only the constants.
\end{proof}

\Needspace{14\baselineskip}
\begin{corollary}\label{cor:dini-family-distance-limits}\index{Metric distance!uniform pseudometric limit}
Let $\omega_0$ be a Hermitian form on $X$ and $d_{\omega_0}$ the distance of the associated Riemannian metric. For each $j$, let $T_j=\theta_j+\ddc v_j\geq0$, where $\theta_j$ is a smooth real $(1,1)$-form, $v_j$ is continuous, and $T_j$ is a smooth Hermitian form off a proper analytic subset $A_j$. Assume
\[
 |\theta_j|_{\omega_0}\leq\Lambda,\qquad
 |v_j(x)-v_j(y)|\leq m(d_{\omega_0}(x,y)),
\]
with one fixed $\Lambda$ and one nondecreasing modulus $m$ tending to zero such that the function in \eqref{eq:dini-distance-modulus} satisfies $D(r_0/2)<+\infty$. By \cref{thm:dini-current-distance}, the identity of $X\setminus A_j$ extends to a continuous surjection of $X$ onto the metric completion of $X\setminus A_j$; let $d_j$ be the pseudometric on $X$ obtained by pulling back the completed distance. Then the diameters of the $d_j$ are uniformly bounded, and a subsequence converges uniformly on $X\times X$ to a continuous pseudometric. If $T_j\geq c\omega_0$ on the regular loci for one $c>0$, the limit is a metric and dominates $\sqrt c\,d_{\omega_0}$.
\end{corollary}
\begin{proof}
A finite fixed chart cover and the theorem give common constants $C,r_1>0$ such that
\[
 d_j(x,y)\leq CD(Cd_{\omega_0}(x,y)),
 \qquad d_{\omega_0}(x,y)<r_1.
\]
For a pair meeting $A_j$, approximate it by regular pairs in the manifold topology. The completed distance and $D$ are continuous, so the same estimate holds for all such pairs in $X\times X$. No uniform bound on the geometry of the analytic sets enters the constants. Dividing a minimizing geodesic of $d_{\omega_0}$ into a uniformly bounded number of short pieces and using the triangle inequality gives a uniform diameter bound. The triangle inequality also gives
\[
 |d_j(x,y)-d_j(x',y')|
 \leq d_j(x,x')+d_j(y,y'),
\]
so the distance functions are equicontinuous. The Arzel\`a--Ascoli theorem gives a subsequence $d_{j_k}$ converging uniformly on $X\times X$ to a function $d_\infty$. Nonnegativity, symmetry, vanishing on the diagonal and the triangle inequality pass to the limit. The lower bound $T_j\geq c\omega_0$ gives $d_j\geq\sqrt c\,d_{\omega_0}$, first for regular pairs and then for all pairs, which proves the last assertion.

Without this lower bound the limit may degenerate. For $\varepsilon_j\downarrow0$, take $\theta_j=\varepsilon_j\omega_0$ and $v_j=0$, with $A_j=\varnothing$. These data satisfy the hypotheses with $m(r)=r^2$, and $d_j=\sqrt{\varepsilon_j}\,d_{\omega_0}$ converges uniformly to zero.
\end{proof}
